\section*{38. Gemeinsam mit R. Brauer und H. Hasse: Beweis eines Hauptsatzes in der Theorie der Algebren}

\begin{center}
\emph{Journal f. d. reine u. angew. Math. 167 (1932), S. 399--404}
\end{center}

\begin{center}
\textbf{Beweis eines Hauptsatzes in der Theorie der Algebren.}\\[0.5em]
Von \emph{R. Brauer} in Königsberg, \emph{H. Hasse} in Marburg und \emph{E. Noether} in Göttingen\footnote{Die Abfassung dieser Note übernahm H. Hasse.}.
\end{center}

Endlich ist es unseren vereinten Bemühungen gelungen, die Richtigkeit des folgenden Satzes zu beweisen, der für die Strukturtheorie der Algebren über algebraischen Zahlkörpern sowie auch darüber hinaus von grundlegender Bedeutung ist:

\paragraph{Hauptsatz.}
\emph{Jede normale Divisionsalgebra über einem algebraischen Zahlkörper ist zyklisch (oder, wie man auch sagt, vom Dicksonschen Typus).}

Es ist uns eine besondere Freude, dieses Ergebnis, als einen im wesentlichen der \(p\)-adischen Methode zu dankenden Erfolg, Herrn Kurt Hensel, dem Begründer dieser Methode, zu seinem 70. Geburtstag vorzulegen.

Unser Beweis besteht in drei Reduktionen, von denen jeder von uns eine beigesteuert hat.\footnote{Diese werden in der Reihenfolge ihrer Entstehung wiedergegeben, die der systematischen Reihenfolge entgegengesetzt ist.}

\paragraph{1.}
Die erste Reduktion gab H. Hasse auf Grund der von ihm kürzlich entwickelten Theorie der zyklischen Algebren über algebraischen Zahlkörpern.\footnote{H. Hasse, \emph{Theorie der zyklischen Algebren über einem algebraischen Zahlkörper}, Gött. Nachr. 1931. Eine ausführliche Darstellung der Beweise, in der insbesondere auch die von E. Noether in einer Vorlesung entwickelte, dort wie hier grundlegende Theorie der Zerfällungskörper und verschränkten Produkte entwickelt wird, erscheint demnächst in den Trans. Amer. Math. Soc. unter dem Titel \emph{Theory of cyclic algebras over an algebraic number field}. Diese Arbeit wird im folgenden mit H zitiert. H, 1--6 machen -- bis auf den Unterschied in der Sprache -- die erstere Note aus.}

\paragraph{Reduktion 1.}
Der Hauptsatz ist bewiesen, wenn gezeigt ist:
\[
\text{I. Jede überall zerfallende Algebra über }\Omega\text{ ist } \sim \Omega .
\]
Zur Abkürzung der Ausdrucksweise soll dabei hier wie im folgenden \glqq{}Algebra \(A\) über \(\Omega\)\grqq{} stets \glqq{}normale einfache Algebra \(A\) über dem algebraischen Zahlkörper \(\Omega\)\grqq{} bedeuten, also ein einfaches hyperkomplexes System mit dem Zentrum \(\Omega\). Ferner bezeichnet \(A\sim\Omega\), daß \(A\) eine vollständige Matrixalgebra in \(\Omega\) ist, also die Zugehörigkeit von \(A\) zu der speziellen, durch \(\Omega\) selbst als Divisionsalgebra über \(\Omega\) bestimmten Klasse im Sinne der Ähnlichkeit (Gleichheit der zugehörigen Divisionsalgebren über \(\Omega\)) [H, 5]. Schließlich ist unter einer überall zerfallenden Algebra über \(\Omega\) eine solche verstanden, bei der für jede Primstelle \(\mathfrak p\) von \(\Omega\) die \(\mathfrak p\)-adische Erweiterung, d.\,h. die durch Erweiterung des Koeffizientenkörpers \(\Omega\) auf den zugehörigen \(\mathfrak p\)-adischen Körper \(\Omega_{\mathfrak p}\) entstehende Algebra, zerfällt:
\[
A_{\mathfrak p}\sim \Omega_{\mathfrak p}.
\]

\paragraph{Beweis.}
Sei \(D\) eine Divisionsalgebra über \(\Omega\). Ist dann \(Z\) ein zyklischer Körper über \(\Omega\) derart, daß für jede Primstelle \(\mathfrak p\) von \(\Omega\) der \(\mathfrak p\)-Grad von \(Z\) ein Multiplum des \(\mathfrak p\)-Index von \(D\) ist [vgl. d. Anm. zu H, 17 Bb], so ist für die zugehörigen Primteiler \(\mathfrak P\) in \(Z\) jeweils der \(\mathfrak P\)-adische Körper \(Z_{\mathfrak P}\) Zerfällungskörper von \(D_{\mathfrak p}\) [H, 18.1]. Daraus ergibt sich, daß die durch Erweiterung des Koeffizientenkörpers von \(D\) auf \(Z\) entstehende Algebra \(D_Z\) über \(Z\) überall zerfällt. Steht nun I bereits fest, für jedes \(\Omega\), so folgt also \(D_Z\sim Z\). Das besagt, daß \(D\) den zyklischen Zerfällungskörper \(Z\) besitzt, also zyklisch darstellbar ist [H, 5]. Dann besitzt \(D\) aber auch einen solchen zyklischen Zerfällungskörper, dessen Grad mit dem Grad von \(D\) selbst übereinstimmt; das heißt, \(D\) ist zyklisch [H, 6, Satz 6]. Unter Voraussetzung I ist der Hauptsatz also bewiesen.

\paragraph{2.}
Für die zweite Reduktion gab R. Brauer den entscheidenden Anstoß, indem er H. Hasse brieflich mitteilte, wie sich die verwandte Frage nach dem genauen Wert des Exponenten einer Divisionsalgebra (die übrigens durch den Hauptsatz mitgelöst wird -- s. u.) mittels des Sylowschen Gruppensatzes auf den Fall eines auflösbaren Zerfällungskörpers zurückführen läßt. Dieser Gedanke ließ sich dann auch zu einer entsprechenden weiteren Reduktion der Zyklizitätsfrage verwenden.

\paragraph{Reduktion 2.}
Satz I ist bewiesen, wenn gezeigt ist:
\[
\text{II. Jede überall zerfallende Algebra mit einem auflösbaren Zerfällungskörper über }\Omega\text{ ist }\sim\Omega .
\]

\paragraph{Beweis.}
Sei \(A\) eine überall zerfallende Algebra über \(\Omega\). Ist dann \(K\) ein galoisscher Zerfällungskörper für \(A\), ferner \(p\) irgendeine Primzahl und \(\Sigma\) einer der auf \(p\) bezüglichen Sylowkörper von \(K\) über \(\Omega\) (d. h. Invariantenkörper einer zu \(p\) gehörigen Sylowgruppe der galoisschen Gruppe), so ist \(A_{\Sigma}\) eine ebenfalls überall zerfallende Algebra über \(\Sigma\), die den auflösbaren Zerfällungskörper \(K\) besitzt. Steht nun II bereits fest (für jedes \(\Omega\)), so folgt also \(A_{\Sigma}\sim \Sigma\). Das besagt, daß \(A\) den Zerfällungskörper \(\Sigma\) von zu \(p\) primem Grade besitzt. Der Index von \(A\) ist daher, als Teiler dieses Grades, zu \(p\) prim. Da dies für jede Primzahl \(p\) gilt, ist er also gleich 1, d. h. \(A\sim\Omega\). Unter der Voraussetzung II ist also I bewiesen.\footnote{Der Gedanke der Reduktion auf auflösbaren Zerfällungskörper mittels des Sylowschen Gruppensatzes wurde schon früher von R. Brauer angewandt, nämlich um zu zeigen, daß jeder Primteiler des Index auch im Exponenten vorkommt (Über den Zusammenhang von arithmetischen und invariantentheoretischen Eigenschaften von Gruppen linearer Substitutionen, Berl. Akad.-Ber. 1926). Neuerdings hat A. A. Albert für diesen Gedanken sowie überhaupt für eine Reihe von allgemeinen Sätzen der R. Brauerschen und E. Noetherschen Theorie einfache von der Darstellungstheorie unabhängige Beweise entwickelt (1. \emph{On direct products, cyclic division algebras, and pure Riemann matrices}; 2. \emph{On direct products}; beides in Trans. Amer. Math. Soc. 33 (1931); für die hier in Rede stehende Reduktion siehe insbesondere Satz 23 in 2.).

\emph{Zusatz während der Drucklegung.} Ferner hat A. A. Albert, im Besitz der brieflichen Mitteilung von H. Hasse, daß der Hauptsatz von ihm für abelsche Algebren bewiesen sei (siehe anschließend im Text), daraus unmittelbar, unabhängig von uns, die folgenden Tatsachen gefolgert:

a) den Hauptsatz für Grade der Form \(2^e\),

b) den unten folgenden Satz 1 (Exponent = Index),

c) neben dem Grundgedanken der Reduktion 2 auch noch den der nachfolgenden Reduktion 3, naturgemäß ohne die Beziehung auf die Reduktion 1, und dementsprechend mit dem Resultat: Für Divisionsalgebren \(D\) vom Primzahlpotenzgrad \(p^e\) über \(\Omega\) gibt es einen Erweiterungskörper \(\Omega'\) von zu \(p\) primem Grade über \(\Omega\), so daß \(D_{\Omega'}\) zyklisch ist.

Alle drei Resultate sind natürlich durch unseren inzwischen geführten Beweis des Hauptsatzes überholt. Sie zeigen jedoch, daß auch A. A. Albert ein unabhängiger Anteil am Beweis des Hauptsatzes zukommt. Schließlich hat A. A. Albert (nach Kenntnis unseres Beweises des Hauptsatzes) noch bemerkt, daß unser zentraler Satz I in ein paar Zeilen aus den Sätzen 13, 10, 9 einer im Druck befindlichen Arbeit von ihm (Bull. Amer. Math. Soc. 37 (1931)) folgt. Der Beweis dieser Sätze beruht im wesentlichen auf denselben Schlüssen, wie unsere Reduktionen 2 und 3.}

\paragraph{3.}
Die dritte Reduktion, die dann, wie H. Hasse -- im Besitz der beiden ersten Reduktionen -- erkannte, zum endgültigen Beweise führt, gab E. Noether, veranlaßt durch eine Mitteilung von H. Hasse; in dieser wurde der Hauptsatz für den Spezialfall eines abelschen Zerfällungskörpers bewiesen, und zwar mittels der allgemeinen Reduktion 1 und formelmäßiger Reduktion des zugehörigen Faktorensystems, als deren Kern E. Noether eben die dritte Reduktion herausschälte. Unabhängig von E. Noether hatte sich übrigens auch R. Brauer diese dritte Reduktion schon überlegt.

\paragraph{Reduktion 3.}
Satz II ist bewiesen, wenn gezeigt ist:
\[
\begin{gathered}\text{III. Jede überall zerfallende Algebra mit einem zyklischen Zerfällungskörper}\\\text{von Primzahlgrad über }\Omega\text{ ist }\sim\Omega .\end{gathered}
\]

\paragraph{Beweis.}
Sei \(A\) eine überall zerfallende Algebra mit einem auflösbaren Zerfällungskörper \(K\) über \(\Omega\). Sei
\[
K=\Lambda_0>\Lambda_1>\cdots>\Lambda_r=\Omega
\]
eine Kette von Körpern zwischen \(K\) und \(\Omega\) derart, daß immer \(\Lambda_i\) über \(\Lambda_{i+1}\) zyklisch von Primzahlgrad ist. Dann ist zunächst \(A_{\Lambda_1}\) eine ebenfalls überall zerfallende Algebra über \(\Lambda_1\), die den zyklischen Zerfällungskörper von Primzahlgrad \(K=\Lambda_0\) besitzt. Steht nun III bereits fest (für jedes \(\Omega\)), so folgt also \(A_{\Lambda_1}\sim \Lambda_1\). Das besagt, daß \(\Lambda_1\) Zerfällungskörper für \(A\) ist. Jetzt beweist man von \(\Lambda_1\) statt \(\Lambda_0\) ausgehend genau so \(A_{\Lambda_2}\sim\Lambda_2\), usw., bis schließlich \(A_{\Lambda_r}\sim\Lambda_r\), d. h. \(A\sim\Omega\). Unter der Voraussetzung III ist also II bewiesen.

\paragraph{4.}
Nun steht aber III nach den Ergebnissen von H. Hasse [H, 24] fest. Also folgt durch die Reduktionen 3, 2, 1 rückwärts die Richtigkeit des Hauptsatzes.
